\documentclass[11pt]{article}
\usepackage[a4paper,margin=25mm]{geometry}
\usepackage{amsmath,amssymb,amsthm}
\usepackage[hidelinks]{hyperref}
\newtheorem{theorem}{Theorem}
\setlength{\emergencystretch}{1em}
\begin{document}
\title{A fourth-order splitting whose negative mass is one half}
\author{Chengzhe Gao}\date{}\maketitle
Conjecture 6557 includes the claim that the sum of negative coefficients
in higher-order splitting is one. We give a counterexample whether this
means the signed sum or the sum of their absolute values.
Order is understood in the usual formal Taylor-series sense for splitting
in noncommuting operators; see \cite{BCM}, Section 3.

Put $S(t)=\exp(tA/2)\exp(tB)\exp(tA/2)$, and define
\[
 T(h)=S(h/6)^4S(-h/3)S(h/6)^4.
\]
Merge consecutive exponentials of $A$. The resulting normalized product
has nineteen alternating, nonzero stages, with the following coefficients:
\[
\begin{array}{c|rrrrrrrrrr}
 A&\frac1{12}&\frac16&\frac16&\frac16&-\frac1{12}&-\frac1{12}
     &\frac16&\frac16&\frac16&\frac1{12}\\[3pt]
 B&\frac16&\frac16&\frac16&\frac16&-\frac13&\frac16&\frac16
     &\frac16&\frac16&
\end{array}
\]
Read the product as $A_1,B_1,A_2,B_2,\ldots,B_9,A_{10}$.
The $A$ coefficients sum to one, as do the $B$ coefficients.

\begin{theorem}
The product $T$ has formal order at least four. Its negative coefficients
have signed sum $-1/2$ and absolute-value sum $1/2$.
\end{theorem}
\begin{proof}
Let $H=A+B$. The symmetric Strang product has formal logarithm
\[
 \log S(t)=tH+t^3C+O(t^5)
\]
for a fixed degree-three Lie polynomial $C$. Indeed its constant,
linear, and quadratic expansion agrees with $\exp(tH)$, and
$S(-t)=S(t)^{-1}$ eliminates all even logarithmic degrees.
If $c_1,\ldots,c_9$ are the nine composition coefficients, then
\[
 \sum c_i=8\cdot\frac16-\frac13=1,\qquad
 \sum c_i^3=8\cdot\frac1{216}-\frac1{27}=0.
\]
Up to degree three the logarithm of the product is therefore
$hH+h^3(\sum c_i^3)C=hH$: commutators of its degree-one terms vanish
because those terms are all multiples of $H$. The composition is
palindromic and hence symmetric, so its logarithm has no degree-four
term either. It follows that $T(h)=\exp(hH)+O(h^5)$.

The only negative coefficients in the merged, alternating product are
$-1/12,-1/3,-1/12$. Their signed sum is $-1/2$, and the sum of their
absolute values is $1/2$. Neither is one.
\end{proof}

This counterexample is consistent with the necessity of some negative
coefficients in higher-order real splitting. It refutes the additional
numerical claim in the source; it does not attempt to refute the usual
necessity theorem. No cancelling zero-time stages or consecutive stages
of the same operator remain in the displayed product.

\paragraph{Formal verification.}
The Lean proof verifies the nineteen-stage product directly and does not
assume a Baker--Campbell--Hausdorff formula. An infinite noncommutative
formal series is a function on all finite words in $A,B$. Exponential
coefficients are $a^n/n!$ on constant-letter words and zero otherwise.
The Cauchy product sums over every prefix--suffix decomposition. Both
reconstruction and completeness of these decompositions are proved.

The code contains all thirty-one words of length at most four and proves
that every arbitrary word below this bound occurs. Rational coefficient
tables certify each successive multiplication. A general prefix theorem
proves that multiplication respects agreement through any given degree;
it connects every intermediate table to the genuine infinite product.
The final table equals $1/n!$ on every word of length $n\leq4$, which is
exactly the coefficient of $\exp(h(A+B))$.

The theorem \texttt{order\_four} therefore proves the universal formal
order condition. A further theorem proves equality of all degree-at-most-
four homogeneous evaluations after arbitrary substitution of word values
in any target algebra. It uses neither commutativity nor special identities
of particular matrices. The component time sums, alternation, nonzero
stages, and both interpretations of the negative-coefficient sum are
checked on the actual coefficient list. Two final theorems negate the
respective signed and absolute-value universal claims.

All certificates are checked by kernel reduction in Lean 4.19.0 using
\texttt{Std} exact rational arithmetic. The accompanying Python program
independently recomputes coefficients with rational fractions; it is not
trusted by Lean. No admitted proofs, new axioms, or external decision
oracle are used.
\begin{thebibliography}{1}
\bibitem{BCM} S. Blanes, F. Casas and A. Murua,
\emph{Splitting and composition methods in the numerical integration
of differential equations}, 2008.
\href{https://arxiv.org/abs/0812.0377}{arXiv:0812.0377}.
\end{thebibliography}
\end{document}
